\section{Introduction}

Autonomous Mobile Robots (AMRs) are rapidly proliferating across service, warehouse logistics, and healthcare automation. A foundational prerequisite for indoor autonomy is local obstacle detection and dynamic avoidance~\cite{liu2024review}. Traditionally, mobile robots rely on active range sensors, predominantly 2D and 3D LiDARs~\cite{li2020lidar}. While LiDAR sensors deliver direct, millimeter-accurate geometric measurements invariant to ambient lighting, their mechanical moving parts, substantial unit cost, physical payload, and electrical power consumption restrict their viability in low-cost, compact robotic platforms.

Conversely, monocular RGB cameras represent an exceptionally lightweight, cost-effective, and energy-efficient sensing modality. However, visual imagery lacks explicit metric depth, rendering direct integration into conventional 2D planar occupancy grid mapping and navigation frameworks non-trivial~\cite{thrun2010probabilistic, elfes1989using}. While RGB-D sensors (e.g., active stereo or structured light cameras) mitigate this gap, their operational range is notoriously limited indoors by optical interference, sunlight dilution, and higher power budgets.

To bridge the interface discrepancy between monocular visual features and planar robot navigation architectures, this paper introduces a real-time \textbf{Vision-Based 2D Scan Generation (V-LiDAR)} system. Instead of estimating dense depth across the entire scene via computationally burdensome deep neural networks, our framework leverages the physical constraint of indoor environments: ground-plane contact. By segmenting traversable floor surfaces and identifying the lowest non-floor pixel in each angular column, obstacle-to-ground contact boundaries are transformed into metric range vectors using precomputed geometric Look-Up Tables (LUTs).

In preliminary work presented as an extended abstract at ICCAS~\cite{lee2025vision_iccas}, we demonstrated the initial feasibility of monocular scan generation. In this expanded article, we substantially advance the state of the art through:
\begin{enumerate}
    \item \textbf{Comprehensive 20-Run Avoidance Verification}: We conduct rigorous multi-trial repeatability evaluations in a 10.0\,m navigation task, improving obstacle avoidance success to 90.0\% (18/20 trials) with an average avoidance onset range of $2.22 \pm 0.08\,\text{m}$.
    \item \textbf{Analytical Geometric Blind Spot Derivation}: We mathematically formalize the near-field ground visibility cutoff ($D_{\text{min}} \approx 2.178\,\text{m}$) as a direct function of mounting height, tilt, and vertical field of view.
    \item \textbf{Costmap Entrapment Analysis}: We analytically dissect why narrow corridors (1.8--2.2\,m) lead to navigation freeze when combined with standard safety inflation radiuses, justifying the experimental decoupling of obstacle detection from hallway wall interference.
\end{enumerate}
